\documentclass[11pt, letterpaper]{article}

\usepackage[top=0.85in, bottom=0.85in, left=0.85in, right=0.85in]{geometry}
\usepackage{amsmath, amssymb, amsthm, mathtools}
\usepackage{longtable}
\usepackage{booktabs}
\usepackage{array}
\usepackage{tabularx}
\usepackage{xcolor}
\usepackage{fancyhdr}
\usepackage{titlesec}
\usepackage{setspace}
\usepackage[T1]{fontenc}
\usepackage{enumitem}
\usepackage{tocloft}
\usepackage{tcolorbox}
\tcbuselibrary{breakable, skins}

% --- Color Definitions ---
\definecolor{RSblue}{RGB}{25, 60, 120}
\definecolor{RSgray}{RGB}{90, 90, 90}
\definecolor{RSgreen}{RGB}{20, 110, 60}
\definecolor{RSred}{RGB}{160, 30, 30}
\definecolor{RSamber}{RGB}{180, 110, 0}
\definecolor{RSpurple}{RGB}{90, 30, 120}

% --- Box Environments ---
\newtcolorbox{rsbox}{
  enhanced, breakable,
  colback=RSblue!4, colframe=RSblue!60,
  title={\textbf{RS Ontological Statement}},
  fonttitle=\small\bfseries\color{white},
  attach boxed title to top left={yshift=-2mm, xshift=4mm},
  boxed title style={colback=RSblue},
  left=4pt, right=4pt, top=6pt, bottom=4pt
}

\newtcolorbox{verifybox}{
  enhanced, breakable,
  colback=RSgreen!4, colframe=RSgreen!50,
  title={\textbf{External Verification (Continuous Lens)}},
  fonttitle=\small\bfseries\color{white},
  attach boxed title to top left={yshift=-2mm, xshift=4mm},
  boxed title style={colback=RSgreen!70!black},
  left=4pt, right=4pt, top=6pt, bottom=4pt
}

\newtcolorbox{observebox}{
  enhanced, breakable,
  colback=RSamber!5, colframe=RSamber!60,
  title={\textbf{Structural Observation}},
  fonttitle=\small\bfseries\color{white},
  attach boxed title to top left={yshift=-2mm, xshift=4mm},
  boxed title style={colback=RSamber!80!black},
  left=4pt, right=4pt, top=6pt, bottom=4pt
}

% --- Page Style ---
\pagestyle{fancy}
\fancyhf{}
\fancyhead[L]{\small\textcolor{RSgray}{RigbySpace Discrete Dynamics}}
\fancyhead[R]{\small\textcolor{RSgray}{Rigby-Tate \& L-Function -- Page \thepage}}
\fancyfoot[L]{\small\textcolor{RSgray}{D. Veneziano}}
\fancyfoot[R]{\small\textcolor{RSgray}{Discrete Unification Framework}}
\renewcommand{\headrulewidth}{0.4pt}

% --- Section Formatting ---
\titleformat{\section}
  {\large\bfseries\color{RSblue}}{\thesection}{1em}{}[\titlerule]
\titleformat{\subsection}
  {\normalsize\bfseries\color{RSblue}}{\thesubsection}{1em}{}
\titleformat{\subsubsection}
  {\normalsize\bfseries\color{RSgray}}{\thesubsubsection}{1em}{}

\onehalfspacing
\setlength{\parskip}{4pt}
\setlength{\parindent}{0pt}

% --- Macros ---
\newcommand{\Z}{\mathbb{Z}}
\newcommand{\SL}{S_{\mathrm{L}}}
\newcommand{\koppa}{\mathbin{\text{\rm\k{o}}}}
\newcommand{\succOp}{\operatorname{succ}}
\newcommand{\PG}{\mathrm{PG}}
\newcommand{\Dcross}{\Delta_{\mathrm{cross}}}
\newcommand{\Tvac}{T_{\mathrm{vac}}}
\newcommand{\Tbary}{T_{\mathrm{bary}}}
\newcommand{\Dg}{D_{\mathrm{gauge}}}
\newcommand{\Sint}{S_{\mathrm{int}}}
\newcommand{\Sigimb}{\Sigma_{\mathrm{imb}}}
\newcommand{\Gmg}{G_{\mathrm{mg}}}
\newcommand{\dslip}{\delta_{\mathrm{slip}}}
\newcommand{\tauX}{\tau_{10}}
\newcommand{\aSC}{\alpha^{-1}_{\mathrm{int}}}
\newcommand{\med}{\oplus}
\newcommand{\aop}{\eta}
\newcommand{\lop}{\lambda}
\newcommand{\psiop}{\psi}
\newcommand{\rank}{\rho}
\newcommand{\Sigs}{\Sigma_{\mathrm{s}}}
\newcommand{\Xis}{\Xi_{\mathrm{s}}}
\newcommand{\Omegas}{\Omega_{\mathrm{s}}}
\newcommand{\Pithree}{\Pi_3}
\newcommand{\nG}{n_G}
\newcommand{\Ksat}{K_{\mathrm{sat}}}
\newcommand{\RD}{R_D}
\newcommand{\bplus}{\boxplus}
\newcommand{\VS}{\mathcal{V}_S}
\newcommand{\piRS}{\pi_{\mathrm{RS}}}

\begin{document}

\begin{titlepage}
  \centering
  \vspace*{2cm}
  {\Huge\bfseries\color{RSblue} RigbySpace Discrete Dynamics\par}
  \vspace{0.5cm}
  {\LARGE The Rigby-Tate Correspondence and the Discrete L-Function\par}
  \vspace{0.5cm}
  {\large A Granular, Forward-Derived Ontological Formalization\par}
  \vspace{2cm}
  \rule{0.6\textwidth}{0.4pt}\par
  \vspace{1cm}
  {\large D.\ Veneziano\par}
  \vspace{0.8cm}
  {\normalsize RigbySpace Discrete Dynamics\par}
  \vspace{1.5cm}
  \rule{0.6\textwidth}{0.4pt}\par
  \vspace{0.5cm}
  {\small $\Tvac = 11 \quad N = 3$ \par}
  \vfill
  {\small\color{RSgray}
  RigbySpace is a formally consistent candidate theory internal to its own
  axioms. No empirical equivalence to external physical theory is asserted
  unless an explicit mapping is provided and derived under that mapping.\par}
\end{titlepage}

\tableofcontents
\newpage

%=====================================================================
\section{Introduction: The Macroscopic Mask of Algebraic Geometry}
%=====================================================================

In standard continuous physics and mathematics, algebraic geometry (specifically the theory of Elliptic Curves) and analytic number theory (L-functions, the Birch and Swinnerton-Dyer conjecture) are treated as fundamental, continuous mathematical truths. The geometry of a circle ($\pi$) is considered an irreducible transcendental constant.

RigbySpace (RS) proves that these continuous, complex mathematical structures are not fundamental. They are the macroscopic, time-averaged masks of a highly efficient, discrete integer engine. The universe does not compute complex scalar multiplication on elliptic curves, nor does it evaluate infinite analytic series. It executes simple integer addition (the Aggregate Transform $\aop$) over a linear temporal dimension.

This document provides the complete, granular, forward-derived formalization of the Rigby-Tate Correspondence and the Discrete L-Function. It demonstrates exactly how the prohibition of Greatest Common Divisor (GCD) reduction forces the accumulation of Structural Friction, how the universe trades operational complexity for temporal duration, how the analytic rank of a curve is deterministically readable from kinematic drift, and how the geometry of a circle ($\pi$) is generated by the vacuum's viscosity.

%=====================================================================
\section{The Constant Lattice Prerequisites}
%=====================================================================

Before deriving the geometric mechanics, we must explicitly construct the required structural constants from the two primitives. Every step uses only integer arithmetic and the imbalance operator $\koppa$ (iterative subtraction).

\begin{rsbox}
\textbf{Primitives:}
\begin{itemize}
    \item Vacuum Period: $\Tvac = 11$.
    \item Triadic Period: $N = 3$.
\end{itemize}

\textbf{1. Mass Gap / Phase Slip ($\Gmg$):}
\[ \Gmg = \dslip = \Tvac \koppa N = 11 \koppa 3 = 2 \]

\textbf{2. 4-4-3 Phase Partition:}
The 11 microticks are partitioned into Emission ($E$), Memory ($M$), and Return ($R$) phases.
\[ N_E = 4, \quad N_M = 4, \quad N_R = 3 \]

\textbf{3. Viscosity Synchronization Window ($\VS$):}
The total structural window over which the three generation levels are encoded.
\[ \VS = \Gmg \times N^2 = 2 \times 3^2 = 2 \times 9 = 18 \]

\textbf{4. Binary Causal Depth ($\rank$):}
For a positive integer $m$, $\rank(m)$ is the unique integer $k \ge 0$ satisfying $2^k \le m < 2^{k+1}$. It is the discrete, non-logarithmic measure of magnitude/complexity, obtained by counting doublings of the mass gap $\Gmg = 2$.
\end{rsbox}

%=====================================================================
\section{Structural Friction and Linear Rank Growth}
%=====================================================================

Because GCD reduction is forbidden, the components of an Explicit Rational Pair (ERP) grow continuously under the action of the Aggregate Transform ($\aop$). This growth is not a mathematical artifact; it is the physical accumulation of causal history.

\begin{definition}[Structural Friction $\Phi(t)$]\label{def:friction}
\begin{rsbox}
The Structural Friction $\Phi(t)$ of the dynamical system at Causal Index $t$ is defined by the maximum Structural Rank of the unreduced integer components of the state $S_t = (n_t, d_t)$:
\[ \Phi(t) := \max(\rank(n_t), \rank(d_t)) \]
This measures the accumulated complexity (causal depth) wedged into the integer components.
\end{rsbox}
\end{definition}

\begin{theorem}[Linear Rank Growth under $\aop$]\label{thm:linear_growth}
\begin{rsbox}
The Structural Rank $\rank(n_t)$ of a state evolving under the Aggregate Transform $\aop(n,d) = (n+d, n)$ grows linearly with the Causal Index $t$. Specifically, $\Phi(t) = \Theta(t)$.

\textbf{Proof:}
Let the initial state be $S_0 = (1, 1)$.
Under $\aop$, the recurrence relation for the numerator is $n_{t+1} = n_t + d_t$.
Since $d_t = n_{t-1}$, this becomes the Fibonacci recurrence:
\[ n_{t+1} = n_t + n_{t-1} \]

\textbf{Upper Bound:}
For $t \ge 2$, $n_{t-1} < n_t$. Therefore:
\[ n_{t+1} < n_t + n_t = 2n_t \]
The Rank function $\rank(n)$ increases by exactly 1 when $n$ doubles. Since $n_{t+1} < 2n_t$, the rank increases by at most 1 per step:
\[ \rank(n_{t+1}) \le \rank(n_t) + 1 \]
Summing over $t$ steps yields the upper bound: $\rank(n_t) < \rank(n_0) + t$.

\textbf{Lower Bound:}
For $t \ge 3$, $n_t = n_{t-1} + n_{t-2}$. Assuming $n_{t-2} > 0$, we have $n_t > n_{t-1}$.
Therefore:
\[ n_{t+1} = n_t + n_{t-1} > n_{t-1} + n_{t-1} = 2n_{t-1} \]
This implies $\rank(n_{t+1}) > \rank(n_{t-1})$. The Rank increases by at least 1 every two steps.
The growth rate is bounded below by $1/2$ per step.

\textbf{Conclusion:}
Since the growth is bounded between $0.5t$ and $1.0t$, the Structural Friction grows strictly linearly:
\[ \Phi(t) = \Theta(t) \]
\end{rsbox}
\end{theorem}

%=====================================================================
\section{The Rigby-Tate Isomorphism}
%=====================================================================

We now establish the exact correspondence between the structural complexity accumulation in RigbySpace and the structural complexity associated with an elliptic curve point in continuous mathematics.

\begin{definition}[Unreduced Height on an Elliptic Curve]\label{def:height}
\begin{rsbox}
For a point $P=(x,y)$ on an elliptic curve, where the coordinate $x=n/d$ is an unreduced rational pair, the Unreduced Height $\hat{h}(P)$ is the maximum Structural Rank of the unreduced integer components:
\[ \hat{h}(P) := \max(\rank(n), \rank(d)) \]
\end{rsbox}
\end{definition}

\begin{lemma}[Quadratic Rank Growth of Scalar Multiplication]\label{lem:quad_growth}
\begin{verifybox}
In standard algebraic geometry, the coordinates of a scalar multiple $mP$ are given by division polynomials. When unreduced, the numerators and denominators of $mP$ are homogeneous polynomials of degree $m^2$ in the coordinates of $P$.
For a homogeneous polynomial of degree $k$ acting on inputs of Rank $R$, the resulting value has a structural Rank proportional to $k \cdot R$.
Therefore, the unreduced height grows quadratically with the scalar $m$:
\[ \hat{h}(mP) = \Theta(m^2) \]
\end{verifybox}
\end{lemma}

\begin{theorem}[The Rigby-Tate Correspondence]\label{thm:rigby_tate}
\begin{rsbox}
The universe does not compute complex scalar multiplication on elliptic curves. It trades operational complexity for temporal duration. The number of linear causal steps $t$ required for the RigbySpace dynamical system to unfold a state complexity equivalent to the structural complexity of $mP$ on an elliptic curve satisfies $t = \Theta(m^2)$.

\textbf{Proof:}
We equate the expressions for Rank growth from the RS substrate (Theorem~\ref{thm:linear_growth}) and the Elliptic Curve structure (Lemma~\ref{lem:quad_growth}).

Let $\Phi(t)$ be the structural friction accumulated over $t$ steps in RigbySpace.
\[ \Phi(t) = C_{\text{dyn}} \cdot t \]

Let $\hat{h}(mP)$ be the unreduced height of the point $mP$.
\[ \hat{h}(mP) = C_{\text{arith}} \cdot m^2 \cdot \hat{h}(P) \]

Equating the two complexity measures:
\[ \Phi(t) = \hat{h}(mP) \]
\[ C_{\text{dyn}} \cdot t = C_{\text{arith}} \cdot m^2 \cdot \hat{h}(P) \]

Solving for $t$:
\[ t = \left( \frac{C_{\text{arith}} \cdot \hat{h}(P)}{C_{\text{dyn}}} \right) \cdot m^2 \]

Since $\hat{h}(P)$, $C_{\text{arith}}$, and $C_{\text{dyn}}$ are constants with respect to $m$, the temporal unfolding $t$ scales exactly quadratically with $m$:
\[ t = \Theta(m^2) \]

\textbf{Conclusion:}
To generate the geometry of $mP$, standard mathematics executes $O(\log m)$ recursive steps of highly complex multiplication. RigbySpace executes $O(m^2)$ steps of simple integer addition ($\aop$). The complex geometry observed in continuous physics is simply the macroscopic, time-averaged shadow of this massive accumulation of linear structural friction.
\end{rsbox}
\end{theorem}

%=====================================================================
\section{The Discrete L-Function ($\mathcal{L}_{\Z}(N)$)}
%=====================================================================

The Birch and Swinnerton-Dyer (BSD) conjecture relates the analytic order of vanishing of a complex L-function to the algebraic rank of an elliptic curve. RigbySpace extracts this exact topological information without real numbers, infinite series, or analytic continuation.

\begin{definition}[Structural Tension and Orientation Sign]\label{def:tension_sign}
\begin{rsbox}
Let the 3D unreduced elliptic state be $S_t = (X_t, Y_t, Z_t)$.
The \textbf{Structural Tension} between consecutive states is the cross-determinant of the Emission ($X$) and Return ($Z$) phases:
\[ \tau_t := X_t Z_{t+1} - X_{t+1} Z_t \in \Z \]

The \textbf{Orientation Sign} $\sigma(m)$ classifies the direction of the cross-determinant turn without invoking absolute values:
\[
\sigma(m) = \begin{cases}
  1, & \text{if } m > 0 \text{ (Left turn)} \\
  0, & \text{if } m = 0 \text{ (Stall)} \\
 -1, & \text{if } m < 0 \text{ (Right turn)}
\end{cases}
\]
\end{rsbox}
\end{definition}

\begin{definition}[The Discrete L-Function]\label{def:discrete_L}
\begin{rsbox}
Fix a positive integer modulus $N$. Let $T_N$ be the $N$-period (the smallest integer $t$ such that the state returns to its initial configuration modulo $N$).
The \textbf{Cumulative Orientation Bias} $R_N$ over one $N$-period is:
\[ R_N := \sum_{t=0}^{T_N-1} \sigma(\tau_t) \in \Z \]

The \textbf{Discrete L-Function} at $N$ is the unreduced Explicit Rational Pair:
\[ \mathcal{L}_{\Z}(N) := (R_N,\; T_N) \in \SL \]
This ERP records the total kinematic drift ($R_N$) against the period length ($T_N$).
\end{rsbox}
\end{definition}

\begin{theorem}[Algebraic Rank from Kinematic Drift]\label{thm:algebraic_rank}
\begin{rsbox}
The algebraic rank of the state's geometry is deterministically readable from the decay profile of $R_N$ relative to $T_N$.

\textbf{1. Rank 0 (Torsion / Closed Loop):}
The trajectory perfectly closes. Every positive tension (left turn) is matched by a negative tension (right turn).
\[ R_N = 0 \implies \mathcal{L}_{\Z}(N) \sim (0, 1) \]
The state exhibits zero kinematic drift.

\textbf{2. Rank 1 (Viscous / Infinite Cyclic):}
The trajectory follows an infinite cyclic subgroup. The signs of $\tau_t$ behave as a bounded random walk. The cumulative bias $R_N$ remains bounded by a constant $C$, while the period $T_N$ grows to infinity as $N$ increases.
\[ |R_N| \le C \text{ while } T_N \to \infty \]
The width between $\mathcal{L}_{\Z}(N)$ and $(0,1)$ decays, confirming asymptotic convergence to zero drift.

\textbf{3. Rank $\ge 2$ (Superfluid / Multiple Generators):}
The existence of multiple independent generators creates perfect symmetry in the lattice. Positive and negative tensions cancel exactly, forcing $R_N = 0$ even at small moduli.
\[ R_N = 0 \implies \mathcal{L}_{\Z}(N) \sim (0, 1) \]

The analytic properties of a system's evolution are strictly encoded in the left/right cross-determinant turns of the unreduced integers.
\end{rsbox}
\end{theorem}

%=====================================================================
\section{The $\pi_{RS}$ Geometry and the Vacuum Viscosity}
%=====================================================================

In continuous mathematics, $\pi$ is an irreducible transcendental constant. In RigbySpace, the geometry of a circle is the macroscopic mask of a specific barycentric oscillation governed by the vacuum's viscosity.

\begin{theorem}[The True Prime-Pair Bounds of $\pi_{RS}$]\label{thm:pi_bounds}
\begin{rsbox}
The standard rational approximant for $\pi$ is $355/113$. However, the ERP $(355, 113)$ is not a frictionless bound because $355$ is not prime ($355 = 5 \times 71$). It is the reduced, macroscopic approximant of a 2:1 barycentric oscillation.

The true frictionless prime-pair bounds ($B_L, B_U$) are derived by solving the Diophantine system for the 2:1 weighted mediant, scaled by the triadic generation depth $N^2 = 9$:
\[ B_L = (1109,\; 353) \]
\[ B_U = (977,\; 311) \]

\textbf{1. Frictionless Verification:}
$1109$ is prime. $353$ is prime.
$977$ is prime. $311$ is prime.
Both $B_L$ and $B_U$ are strictly frictionless prime-pair ERPs.

\textbf{2. The 2:1 Barycentric Average:}
The phase indicator $\phi = C_R \koppa 2$ enforces a 2:1 occupation ratio favoring the lower bound.
Numerator: $2(1109) + 977 = 2218 + 977 = 3195$.
Denominator: $2(353) + 311 = 706 + 311 = 1017$.
The unreduced barycentric core is $(3195, 1017)$.
Dividing by the triadic scaling factor $N^2 = 9$:
$3195 / 9 = 355$.
$1017 / 9 = 113$.
This exactly recovers the macroscopic approximant $(355, 113)$.

\textbf{3. The Bracket Capacity and Vacuum Viscosity:}
The structural tension (Bracket Capacity) between the true prime-pair bounds is measured by their cross-determinant:
\[ \Dcross(B_U, B_L) = (977 \times 353) - (1109 \times 311) \]
\[ \Dcross(B_U, B_L) = 344881 - 344899 = -18 \]

The magnitude of this structural tension is exactly $18$.
In the RS constant lattice, $18 = \VS$ (The Viscosity Synchronization Window).

\textbf{Conclusion:}
The geometry of a circle ($\pi$) is not an independent mathematical axiom. It is the direct physical byproduct of the vacuum's viscosity ($\VS = 18$) forcing a 2:1 barycentric oscillation between the prime-pair bounds $(1109, 353)$ and $(977, 311)$.
\end{rsbox}
\end{theorem}

\end{document}
